% !TeX root = ../papers/nc-p3.tex
\section{被約な低次数の点曲線と三角形の全持上げ}
\label{sec:reduced-cubic-point-curves}
前節と同じ点ファイバー条件のもとで、点グラフが純な被約曲線 $Y$ と交わらない有限部 $F$ の和とする。

\begin{theorem}\label{thm:reduced-cubic-point-curves}
$1\le\deg Y\le3$ では、滑らかな平面三次の既出の枝以外に、三角形の次の枝だけがある。
$k^*\setminus\{1\}$ の元 $\lambda$ に対して $q^3=\lambda$ を一つ選び、添字を法3として
\[
              r_i=x_{i+1}x_{i+2}-q x_i^2
\]
を三生成商の関係とする（向きは頂点座標の選択による）。
各 $q$ について中心文字 $t$ を加え、
\[
       r_i=t\sum_j m_{ij}x_j+t^2 b_i,\qquad m_{ij}=m_{ji},\quad m_{ii}=0
\]
とする全六係数が採用される。いずれも裸周期1、全次数で標準Koszul左右AS正則であり、$F$ の長さは8である。
異なる $\lambda$ の間に添字保持同型はない。固定 $\lambda$ の同型類は分裂点と
$\mathbb P(1,1,1,2,2,2)/(\mathbb Z/3)^2$ である。群作用は文字の巡回置換と
$\operatorname{diag}(1,\zeta,\zeta^2)$（$\zeta$ は原始三乗根）の関係への代入で与える。
この枝の全点は滑らかな平面三次の枝からの平坦特殊化である。
\end{theorem}
\begin{proof}[曲線の被覆]
線と滑らかな二次曲線の開閉成分は系\ref{cor:smooth-point-curves-covered}の証明で除かれる。
節点を高々一本しか持たない滑らかな有理成分（正規化上で枝を数える）も
系\ref{cor:nodal-fixed-count}の証明をその成分に適用すれば、次数が足りない。
従って二直線、平面を異にする二次曲線と直線、三直線の鎖などの有理な端を持つ配置は現れない。
整な非平面三次は有理正規三次である（正規化の次数3の束の独立四切断は種数0の完全系を与える）。
三直線が一点で交わり平面を張らない場合の局所式は $xy=xz=yz=0$ であり、
各枝上の余接層は自由部分と長さ2の捩れを持つ。成分上の恒等は定理\ref{thm:fixed-differential}の第二項に反し、
それ以外では第一項の固定長は零だが唯一の三重点は固定される。これで非平面の場合も尽きる。

被約な特異平面三次では $\omega_Y=\mathcal O_Y$ であり、全成分の正規化は有理である。
局所完全交叉ゆえ各正規化上で $M(T)=\nu^*\omega_Y$ である。実際、余法束列を正規化に引き戻すと
余法束の射は一般点で単射であり、自由な始域上で単射となるから、行列式は随伴公式の束を返す。
点ごとに恒等となる成分があれば固定微分の第二式で $\deg M=0$ を得るが $M\subset\omega_{\mathbb P^1}$ で矛盾する。
従って固定スキームは次数零の有効因子、従って空である。唯一の特異点を持つ節点・尖点の整三次、
接する二次曲線と直線、共点の平面三直線は除かれる。横断する二次曲線と直線では次数の異なる各成分を保つが、
射影直線の自己同型は固定点を持つ。残る三角形では四移動は三頂点と三辺を巡回置換する。
三直線が相異なる三節点で会えば全位置で同一平面に入る。

\smallskip
\noindent\textit{平面を消す文字の正規性。}
三角形の頂点を点列に沿って標識し、各辺の両端の文字を対応する座標とする。
隣接平面に制限した双一次関係は三次元以下である。各頂点で対角係数が消え、
二頂点間の各辺ではその二つの交差係数の間に一つの比例式があるためである。
二式以下なら平面の積で曲面が残るので、制限はこの三次元全体である。基底は
$x_ax'_b-c_{ab}x_bx'_a$ ($a<b$, $c_{ab}\ne0$) となる。
左右行列式は三座標の積の非零倍となる（恒等零なら平面全体が点射影像に入る）。
各曲線点で階数2ゆえ、平面の積での零スキームは三角形のグラフである。

平面を消す文字を $t_i$ とする。制限で消える三関係を定理\ref{thm:plane-cubic-point-graphs}の
$A_i,B_i$ の表示で書く。$A_i$ の核に平面内の点 $P$ があるとする。
$P$ が三角形の滑らかな点なら、三混合式の微分は次の点との対で階数1以下となり、
全六式の微分階数が5に達しない。核に頂点があれば、標準四移動が平面も混合関係も保つため
三頂点が全て核に入る。すると $A_i=0$ となり滑らかな点でも同じ矛盾を得る。
ここで核は $t_i=0$ にある前点で全混合式が後ろの任意の文字に対して消える条件なので内在的である。
$P$ が三角形外なら次点と前点は共に平面外、かつ $B_{i-1}P=0$。
前掲定理の純な三次交叉の矛盾によりこれも不可能である。逆向きも同様で両行列は可逆となる。

従って $t$ の両側イデアルは片側から生成され、三文字商の三次元数は $20-10=10$ 以上である。
先ほどの三つの交換関係の三次交叉の次元は高々1である。
実際、両方の隣接対で文字の交換は非零倍率を与え、隣接文字が一致する語の係数は零である。
文字が重複する全語で零、相異なる三文字の六語の係数は一つから全て決まる。
よって次元はちょうど1、三次元数は10、三次交叉は左右で全支持となる。
$t$ の二次から三次への左右単射性を得たので、前掲定理と同じ二回の移動で通常の中心表示に戻る。
補題\ref{lem:ternary-finite-graph-test}と定理\ref{thm:regular-normal-linear}の三次PBW条件により
三文字商と四文字の持上げは全次数で正則である。

\smallskip
\noindent\textit{中心表示での移動と同型。}
中心表示の三文字商の移動 $\rho$ は三角形上で固定点を持たない。
これには節点での接空間を仮定しない。曲線のグラフは平面の積で三式の完全交叉であり、
周囲の二平面方程式を加えて局所完全交叉である。従って六式の超過交叉の寄与は
$\deg(\omega_Y^{-1}(2,2))=12$ である（余法束の行列式または系\ref{cor:point-graph-excess}の法束計算）。
よって孤立部の長さは8である。中心関係によりこれらは全てアフィン対角 $t=1$ 上にあり、
そこで三つの二次式の孤立共通零スキームになる。対角の射影閉包で三式が無限遠にも零点を持てば
アフィン孤立部の長さは8未満となる。実際、一般の二つの線型結合は共通曲面を持たず
（三式の共通曲面は有限なアフィン部にも無限遠の三角形にも入らない）、
次数4の完全交叉曲線を与える。共通曲線成分があればその正次数分を除いて第三式と交わる
残りのサイクルの交数は8未満であり、各アフィン孤立長を含む。共通曲線がなければ
無限遠点の正の交数を8から引くことになる。よって $\rho$ は固定点を持たず、三辺を巡回する。

頂点の順と座標を選び、$\rho$ が頂点を一つ送る規約にそろえる。
辺の座標比の倍率から関係は $x_{i+1}x_{i+2}-q_i x_i^2$ となる
（頂点の巡回の向きは表示に合わせて取る）。$q_i$ は全て非零、積 $\lambda$ が1なら
行列式が恒等零である。対角座標変更で $q_i$ は $d_{i+1}d_{i+2}/d_i^2$ の倍率を受け、
その像は積1の全二次元トーラスである。従って三つとも $q^3=\lambda$ の同じ $q$ にできる。
この積は同型不変である。平面の同型は三角形を保つ単項行列であり、
標準の巡回を標準の巡回に運ぶ頂点置換は巡回置換だけである。同じ $q$ を保つ対角変換は
共通倍率を除き $\operatorname{diag}(1,\zeta,\zeta^2)$ のべきで尽きる。

逆に $q^3\ne1$ では三式の左右行列式は三角形を簡約に切り、各点で階数2である。
巡回syzygy $\sum x_i r_i=\sum r_i x_i$ は両端で全支持なので三生成商は正則である。
系\ref{cor:sklyanin-central-classification}の係数計算は、純平方係数 $-q$ が非零で
巡回syzygyを持つ今の三式にも適用できる。対称行列と任意の定数欄を得て、
文字への $t$ の加算で対角を一意に零にできる。固定点のない平面グラフの近くでは
中心関係が両位置の $t$ を消し、対角の三式は無限遠に零点を持たず長さ8の完全交叉となる。
三角形の張る平面は内在的なので、線を指定しない全ての同型についても上の有限群と
重み $(1,2)$ の共通倍率で尽きる。

最後に $r_i$ に $u x_{i+2}x_{i+1}$ を加える。三次の行列式は全体倍率を除き
$(1+u^3-q^3)x_0x_1x_2+uq(x_0^3+x_1^3+x_2^3)$ である（$u$ の符号規約による全体倍率を許す）。
一般の $u$ で滑らかなHesse三次となり、対角制限も一般に固定点を持たない
（存在すれば $(1+u)^3=q^3$）。三次syzygyと同じ六補正係数は有効で、
一般点と $u=0$ の全次数のHilbert関数が一定である。各次数の有限表示加群は
$k[u]_{(u)}$ 上で自由なので、指定した全点への平坦特殊化を得る。
\end{proof}
